\documentclass[10pt,a4paper]{amsart}
\usepackage{amsmath,amssymb,amsthm,graphicx}
% Retained minimal preamble for the Stacks Project chapter varieties.tex.
% Verbatim excerpt of the chapter preamble: Xy-pic setup (native preamble lines 14-18), the native theorem
% declarations with the mathematical macros (native preamble lines 154-217), and the project own no-op
% definitions of the reference/slogan/history display environments (native wrapper.tex lines 1-3) that the
% counted statement uses for its slogan. Copied mechanically; nothing added, removed or redefined.
% For commutative diagrams we use Xy-pic
\usepackage[all]{xy}

% We use 2cell for 2-commutative diagrams.
\xyoption{2cell}
% Theorem environments.
%
\theoremstyle{plain}
\newtheorem{theorem}[subsection]{Theorem}
\newtheorem{proposition}[subsection]{Proposition}
\newtheorem{lemma}[subsection]{Lemma}

\theoremstyle{definition}
\newtheorem{definition}[subsection]{Definition}
\newtheorem{example}[subsection]{Example}
\newtheorem{exercise}[subsection]{Exercise}
\newtheorem{situation}[subsection]{Situation}

\theoremstyle{remark}
\newtheorem{remark}[subsection]{Remark}
\newtheorem{remarks}[subsection]{Remarks}

\numberwithin{equation}{subsection}

% Macros
%
\def\lim{\mathop{\mathrm{lim}}\nolimits}
\def\colim{\mathop{\mathrm{colim}}\nolimits}
\def\Spec{\mathop{\mathrm{Spec}}}
\def\Hom{\mathop{\mathrm{Hom}}\nolimits}
\def\Ext{\mathop{\mathrm{Ext}}\nolimits}
\def\SheafHom{\mathop{\mathcal{H}\!\mathit{om}}\nolimits}
\def\SheafExt{\mathop{\mathcal{E}\!\mathit{xt}}\nolimits}
\def\Sch{\mathit{Sch}}
\def\Mor{\mathop{\mathrm{Mor}}\nolimits}
\def\Ob{\mathop{\mathrm{Ob}}\nolimits}
\def\Sh{\mathop{\mathit{Sh}}\nolimits}
\def\NL{\mathop{N\!L}\nolimits}
\def\CH{\mathop{\mathrm{CH}}\nolimits}
\def\proetale{{pro\text{-}\acute{e}tale}}
\def\etale{{\acute{e}tale}}
\def\QCoh{\mathit{QCoh}}
\def\Ker{\mathop{\mathrm{Ker}}}
\def\Im{\mathop{\mathrm{Im}}}
\def\Coker{\mathop{\mathrm{Coker}}}
\def\Coim{\mathop{\mathrm{Coim}}}

% Boxtimes
%
\DeclareMathSymbol{\boxtimes}{\mathbin}{AMSa}{"02}

%
% Macros for moduli stacks/spaces
%
\def\QCohstack{\mathcal{QC}\!\mathit{oh}}
\def\Cohstack{\mathcal{C}\!\mathit{oh}}
\def\Spacesstack{\mathcal{S}\!\mathit{paces}}
\def\Quotfunctor{\mathrm{Quot}}
\def\Hilbfunctor{\mathrm{Hilb}}
\def\Curvesstack{\mathcal{C}\!\mathit{urves}}
\def\Polarizedstack{\mathcal{P}\!\mathit{olarized}}
\def\Complexesstack{\mathcal{C}\!\mathit{omplexes}}
% \Pic is the operator that assigns to X its picard group, usage \Pic(X)
% \Picardstack_{X/B} denotes the Picard stack of X over B
% \Picardfunctor_{X/B} denotes the Picard functor of X over B
\def\Pic{\mathop{\mathrm{Pic}}\nolimits}
\def\Picardstack{\mathcal{P}\!\mathit{ic}}
\def\Picardfunctor{\mathrm{Pic}}
\def\Deformationcategory{\mathcal{D}\!\mathit{ef}}
\newenvironment{reference}{}{}
\newenvironment{slogan}{}{}
\newenvironment{history}{}{}
\usepackage{cleveref}
\crefname{originallemma}{original lemma}{original lemmas}
\Crefname{originallemma}{Original lemma}{Original lemmas}
\makeatletter
\newlabel{algebra-lemma-characterize-separable-field-extensions}{{030W}{source}{}{sourcepointer-1}{}}
\newlabel{algebra-lemma-characterize-separable-field-extensions@cref}{{[originallemma][0][]030W}{[1][0][]source}{}{}{}}
\makeatother
\crefname{originallemma}{original lemma}{original lemmas}
\Crefname{originallemma}{Original lemma}{Original lemmas}
\makeatletter
\newlabel{algebra-lemma-separable-smooth}{{00TV}{source}{}{sourcepointer-2}{}}
\newlabel{algebra-lemma-separable-smooth@cref}{{[originallemma][0][]00TV}{[1][0][]source}{}{}{}}
\makeatother
\crefname{originallemma}{original lemma}{original lemmas}
\Crefname{originallemma}{Original lemma}{Original lemmas}
\makeatletter
\newlabel{lemma-generic-points-geometrically-reduced}{{04KS}{source}{}{sourcepointer-3}{}}
\newlabel{lemma-generic-points-geometrically-reduced@cref}{{[originallemma][0][]04KS}{[1][0][]source}{}{}{}}
\makeatother
\crefname{originallemma}{original lemma}{original lemmas}
\Crefname{originallemma}{Original lemma}{Original lemmas}
\makeatletter
\newlabel{lemma-smooth-geometrically-normal}{{056T}{source}{}{sourcepointer-4}{}}
\newlabel{lemma-smooth-geometrically-normal@cref}{{[originallemma][0][]056T}{[1][0][]source}{}{}{}}
\makeatother
\crefname{originallemma}{original lemma}{original lemmas}
\Crefname{originallemma}{Original lemma}{Original lemmas}
\makeatletter
\newlabel{properties-lemma-characterize-integral}{{01ON}{source}{}{sourcepointer-5}{}}
\newlabel{properties-lemma-characterize-integral@cref}{{[originallemma][0][]01ON}{[1][0][]source}{}{}{}}
\makeatother
\title{A reduced locally finite type scheme over a field is geometrically reduced exactly when it contains a dense smooth open}
\author{The Stacks Project Authors}
\date{stacks-project:varieties.tex:local-snapshot-20261011; GNU FDL 1.2 or later}
\begin{document}
\maketitle
\noindent\textit{Editorial note.} All delivered mathematical text is the author's own, copied from the pinned Stacks Project chapter file varieties.tex; no statement or proof step was rewritten, reordered or removed.\par\medskip
\noindent\textit{Editorial note.} The proof is complete as written and ends with the existence of a smooth principal localization; no step is referred to the reader.\par\medskip
\noindent\textit{Editorial note.} The counted spans carry 5 navigation stubs and no bibliography citation.\par\medskip
\noindent\textit{Editorial note.} The preamble is the chapter's retained minimal preamble (sources/stacks-varieties/retained-preamble.tex lines 1-77), a verbatim excerpt of the native preamble.\par\medskip
\noindent\textit{Editorial note.} The label lemma-geometrically-reduced-dense-smooth-open is defined only in varieties.tex, so no cross-chapter restatement had to be weighed for this card.\par\medskip
\noindent\textit{Allowed standard prerequisites (editorial).}\begin{itemize}
\item Smoothness implies geometric normality and geometric reducedness. Same chapter, Lemma lemma-smooth-geometrically-normal, invoked at proof lines 4704-4705. Supplies that the dense smooth open is geometrically reduced at the generic points.
\item Geometric reducedness can be tested at generic points. Same chapter, Lemma lemma-generic-points-geometrically-reduced, invoked at proof lines 4706-4707. Supplies the passage from generic points to the whole scheme.
\item A geometrically reduced field extension of a field is separable. Algebra, Lemma algebra-lemma-characterize-separable-field-extensions, invoked at proof lines 4719-4720. Supplies separability of the function field over k.
\item Separable residue field implies smoothness at a prime for finite type algebras. Algebra, Lemma algebra-lemma-separable-smooth, invoked at proof lines 4722-4724. Supplies that the affine chart is smooth at the zero ideal, hence that a principal localization is smooth.
\end{itemize}\medskip
\ifdefined\hypertarget\hypertarget{sourcepointer-1}{}\fi
\noindent\textit{Source reference.} Original lemma 030W: sources/stacks-chow/tags.txt, source lines 3963--3963. Pointer outside the counted span, invoked once for context; the load-bearing argument of the counted statement is the author's proof reproduced in full in the proof span.\par\medskip
\ifdefined\hypertarget\hypertarget{sourcepointer-2}{}\fi
\noindent\textit{Source reference.} Original lemma 00TV: sources/stacks-chow/tags.txt, source lines 1089--1089. Pointer outside the counted span, invoked once for context; the load-bearing argument of the counted statement is the author's proof reproduced in full in the proof span.\par\medskip
\ifdefined\hypertarget\hypertarget{sourcepointer-3}{}\fi
\noindent\textit{Source reference.} Original lemma 04KS: sources/stacks-varieties/pinned-original-source.tex, source lines 555--555. Pointer outside the counted span, invoked once for context; the load-bearing argument of the counted statement is the author's proof reproduced in full in the proof span.\par\medskip
\ifdefined\hypertarget\hypertarget{sourcepointer-4}{}\fi
\noindent\textit{Source reference.} Original lemma 056T: sources/stacks-varieties/pinned-original-source.tex, source lines 4617--4617. Pointer outside the counted span, invoked once for context; the load-bearing argument of the counted statement is the author's proof reproduced in full in the proof span.\par\medskip
\ifdefined\hypertarget\hypertarget{sourcepointer-5}{}\fi
\noindent\textit{Source reference.} Original lemma 01ON: sources/stacks-chow/tags.txt, source lines 2200--2200. Pointer outside the counted span, invoked once for context; the load-bearing argument of the counted statement is the author's proof reproduced in full in the proof span.\par\medskip
\begin{lemma}
\label{lemma-geometrically-reduced-dense-smooth-open}
Let $X$ be a reduced scheme that is locally of finite type over a field $k$.
Then $X$ is geometrically reduced over $k$ if and only if
$X$ contains a dense open which is smooth over $k$.
\end{lemma}

\begin{proof}
The problem is local on $X$, hence we may assume $X$ is quasi-compact.
Let $X = X_1 \cup \ldots \cup X_n$ be the irreducible components of $X$.
First suppose that $X$ contains a dense open which is smooth over $k$.
Then $X$ is geometrically reduced at the generic point of each $X_i$
(for example by Lemma \ref{lemma-smooth-geometrically-normal}).
By Lemma \ref{lemma-generic-points-geometrically-reduced},
it follows that $X$ is geometrically reduced.

\medskip\noindent
Conversely, suppose that $X$ is geometrically reduced.
Since $Z = \bigcup_{i \not = j} X_i \cap X_j$ is nowhere dense in $X$,
we may replace $X$ by $X \setminus Z$. As $X \setminus Z$ is a
disjoint union of irreducible schemes, this reduces us to the case
where $X$ is irreducible. As $X$ is irreducible and reduced, it is
integral, see
Properties, Lemma \ref{properties-lemma-characterize-integral}.
Let $\eta \in X$ be its generic point.
Then the function field $K = k(X) = \kappa(\eta) = \mathcal{O}_{X, \eta}$
is geometrically reduced over $k$, hence separable over $k$, see
Algebra, Lemma \ref{algebra-lemma-characterize-separable-field-extensions}.
Let $U = \Spec(A) \subset X$ be any nonempty affine open
so that $K = A_{(0)}$ is the fraction field of $A$. Apply
Algebra, Lemma \ref{algebra-lemma-separable-smooth}
to conclude that $A$ is smooth at $(0)$ over $k$.
By definition this means that some principal localization
of $A$ is smooth over $k$ and we win.
\end{proof}

\end{document}
